% Complete Chapter 5 introduction, source PDF 61 / printed 51.
% block: approximation-hilbert-heading page:61 type:chapter
\clearpage\setcounter{chapter}{4}
\chapter{{Hilbert 复形的逼近}}\label{chap:approximation-hilbert-complexes}

% block: approximation-hilbert-goal page:61 type:ordinary-paragraph
{现在的目标是设计和分析数值方法, 以近似求解与某个闭 Hilbert 复形相伴的 Hodge Laplacian 问题. 我们将采用 Galerkin 方法, 即用有限维子空间替换原始或混合弱表述中出现的试探空间和检验空间. 在实际应用中, 当 Hilbert 复形为 de Rham 复形或类似复形时, 所考虑的子空间将是有限元空间.}

% block: approximation-short-complex page:61 type:ordinary-paragraph
{要定义 Hilbert 复形在某个层次 $k$ 上的 Hodge Laplace 问题, 只需要从 $W^{k-1}$ 到 $W^k$, 定义域为 $V^{k-1}$ 的算子 $d^{k-1}$, 以及从 $W^k$ 到 $W^{k+1}$, 定义域为 $V^k$ 的算子 $d^k$; 参见 \eqref{eq:abstract-hodge-operator}. 因此, 为避免引入不必要的负担, 本章不需要整个 Hilbert 复形 \eqref{eq:hilbert-domain-complex}, 而只需要其中的一个短片段:}
\begin{equation}\label{eq:approximation-short-complex}
V^{k-1}\xrightarrow{d}V^k\xrightarrow{d}V^{k+1}.
\end{equation}

% block: approximation-chapter-outline page:61 type:ordinary-paragraph
{我们将看到, 原始表述的 Galerkin 离散化存在问题, 因此将着重讨论混合表述的 Galerkin 离散化. 这需要两个相容的 Galerkin 子空间, 分别对应 $V^{k-1}$ 和 $V^k$. 本章的主要结果是对子空间的选取建立假设, 保证 Galerkin 方法的相容性和稳定性, 并由此导出推论. Box~\ref{box:approximation-hilbert-summary} 概括了本章的主要结果.}
